A number of future modifications suggested themselves through this project.

Modifying the subgradient algorithm, for instance by using the r-algorithm as is done in 
\cite{wu2006capacitated} is one possibility, although it is unlikely this would drastically
improve the LR ESLP such that the method becomes comparable to the gurobi solver.

A more interesting modification would be to introduce time-dependency on when sites can be setup, 
rather than assuming all sites are constructed at initiation. This would make such modelling far more
relevant for capital-cost forecasting, as prolonging a site's 
construction until such a time as it is needed could then alter its setup-costs through
time-value-of-money considerations. This could result in more realistic solutions being optimal.
For instance, if demand were expected to rise at some future date at one site drastically, delaying
constructiong to meet the additional demand until it is realised could result in lower
total cost.

A very practical improvement would be to extend this model to consider storage capabilities at different
sites too, incorporating different storage technologies and associated costs and savings in 
utilising them. This could potentially make some facilities like solar and wind more appealing,
although the additional complexity of considering at every site an optimal storage technology
would significantly increase the size of the problem.
